%----------------------------------------------------------------------------------------
% NOTATIE EN FIGUREN VAN DE CURSUS RUIMTEMEETKUNDE (2de graad)
% Staat in de preambule en niet in een hoofdstuk: ook de lijsten met oplossingen gebruiken dit.
% Vertrekpunt: notatie.tex van Driehoeksmeting (2de graad), met figuren in cavalièreperspectief erbij.
%----------------------------------------------------------------------------------------
% Bewaarde grafiek, geplaatst door maak-pdf.py: bestand en pagina in figuren-cache/
\newcommand{\grafiekbestand}[2]{\includegraphics[page=#2,trim=3bp 3bp 3bp 3bp]{figuren-cache/#1}}
\usetikzlibrary{angles,quotes}

% Lengte van een lijnstuk: \lengte{AB} geeft |AB|; \abs{-2} de absolute waarde.
\newcommand{\lengte}[1]{\lvert #1\rvert}
\newcommand{\abs}[1]{\lvert #1\rvert}
% Hoek met een dakje: \hoek{A}
\newcommand{\hoek}[1]{\widehat{#1}}

% Zoals stappen, maar met een bredere kolom voor de berekening: voor lange formules
\newenvironment{stappenbreed}{%
  \par\renewcommand{\arraystretch}{1.25}%
  \noindent\begin{tabular}{@{}>{$\displaystyle}p{9.8cm}<{$}@{\hspace{4mm}}>{\RaggedRight\small\itshape\arraybackslash}p{4.9cm}@{}}%
}{\end{tabular}\par}

% Zoals stappen, met eigen kolombreedtes: \begin{stappenvrij}{6.2cm}{3.8cm}. Voor naast een figuur.
\newenvironment{stappenvrij}[2]{%
  \par\renewcommand{\arraystretch}{1.25}%
  \noindent\begin{tabular}{@{}>{$\displaystyle}p{#1}<{$}@{\hspace{3mm}}>{\RaggedRight\small\itshape\arraybackslash}p{#2}@{}}%
}{\end{tabular}\par}

% Redeneerstappen naast een figuur: tekst in de linkerkolom (geen formules), uitleg cursief rechts.
% \begin{redenering}{6cm}{4cm} stap & reden \\ ... \end{redenering}
\newenvironment{redenering}[2]{%
  \par\renewcommand{\arraystretch}{1.2}%
  \noindent\begin{tabular}{@{}>{\RaggedRight\arraybackslash}p{#1}@{\hspace{3mm}}>{\RaggedRight\small\itshape\arraybackslash}p{#2}@{}}%
}{\end{tabular}\par}

% Tekst of berekening links, figuur rechts: \naast{links}{rechts}
\newsavebox{\naastbox}
\newcommand{\naast}[3][]{%
  \par\noindent\sbox{\naastbox}{#3}%
  \begin{minipage}[t]{\dimexpr\linewidth-\wd\naastbox-3mm\relax}\vspace{0pt}#2\end{minipage}\hfill
  \begin{minipage}[t]{\wd\naastbox}\vspace{0pt}\usebox{\naastbox}\end{minipage}\par\addvspace{1mm}}
% Een rij figuren, gecentreerd
\newcommand{\figuren}[1]{\par\addvspace{1mm}{\centering #1\par}\addvspace{1.5mm}}
% Opgave met een figuur rechts: \begin{metgrafiek}[breedte]{tikz} opgave en deelvragen \end{metgrafiek}
\newsavebox{\grafiekbox}
\newenvironment{metgrafiek}[2][]{%
  \sbox{\grafiekbox}{#2}%
  \begin{minipage}[t]{\dimexpr\linewidth-\wd\grafiekbox-3mm\relax}%
}{%
  \end{minipage}\hfill
  \begin{minipage}[t]{\wd\grafiekbox}\vspace{-2mm}\usebox{\grafiekbox}\end{minipage}\par}

% Een punt als stip: \stip{x}{y}, of met label: \punt[positie]{x}{y}{label}
\newcommand{\stip}[2]{\fill (#1,#2) circle (1.5pt);}
\newcommand{\punt}[4][above]{\fill (#2,#3) circle (1.5pt) node[#1,inner sep=2pt] {$#4$};}
% Een benoemd punt (coördinaat) als stip met label: \pnt[positie]{P}{$P$}
\newcommand{\pnt}[3][above]{\fill (#2) circle (1.6pt) node[#1,inner sep=1.5pt,text=black] {#3};}

%----------------------------------------------------------------------------------------
% RUIMTELIJKE FIGUREN IN CAVALIÈREPERSPECTIEF
% Optie [schuin] bij tikzpicture: x naar rechts, y omhoog, z naar achter. De wijkende richting z staat
% onder 45 graden en is half zo lang (0,5 . cos 45 = 0,3536).
%----------------------------------------------------------------------------------------
\tikzset{schuin/.style={x={(1cm,0cm)},y={(0cm,1cm)},z={(.3536cm,.3536cm)}},
  ribbe/.style={grijs,semithick,line join=round},
  verborgen/.style={lijngrijs,dashed},
  nadruk/.style={accent,line width=1.6pt,line cap=round},
  snede/.style={fill=accent!30,fill opacity=.8,draw=accent,thick,line join=round},
  vlakdeel/.style={fill=accent,fill opacity=.2},
  vlakvulling/.style={fill=accent!10,draw=accent!60,line join=round}}
% De hoekpunten van een balk ABCD.EFGH als coördinaten (A) ... (H): breed a, hoog b, diep c.
% ABCD is het grondvlak (A vooraan links, B vooraan rechts, C achteraan rechts, D achteraan links),
% E ligt boven A, F boven B, G boven C en H boven D. D is het verborgen hoekpunt.
\newcommand{\balkpunten}[3]{%
  \coordinate (A) at (0,0,0); \coordinate (B) at (#1,0,0); \coordinate (C) at (#1,0,#3); \coordinate (D) at (0,0,#3);
  \coordinate (E) at (0,#2,0); \coordinate (F) at (#1,#2,0); \coordinate (G) at (#1,#2,#3); \coordinate (H) at (0,#2,#3);}
% Een doorzichtige balk met de verborgen ribben gestippeld: \draadbalk{a}{b}{c}; zet ook (A) ... (H).
\newcommand{\draadbalk}[3]{%
  \balkpunten{#1}{#2}{#3}%
  \fill[accent!5] (A) -- (B) -- (C) -- (G) -- (H) -- (E) -- cycle;
  \draw[verborgen] (D) -- (A) (D) -- (C) (D) -- (H);
  \draw[ribbe] (A) -- (B) -- (F) -- (E) -- cycle (B) -- (C) -- (G) -- (F) (E) -- (H) -- (G);}
% Een kubus met ribbe s: \draadkubus{s}
\newcommand{\draadkubus}[1]{\draadbalk{#1}{#1}{#1}}
% De namen A ... H bij de hoekpunten
\newcommand{\balknamen}{%
  \node[below left,inner sep=1.5pt] at (A) {$A$}; \node[below right,inner sep=1.5pt] at (B) {$B$};
  \node[right,inner sep=2pt] at (C) {$C$}; \node[left,inner sep=1.5pt] at (D) {$D$};
  \node[left,inner sep=2pt] at (E) {$E$}; \node[below right,inner sep=1.5pt] at (F) {$F$};
  \node[above right,inner sep=1.5pt] at (G) {$G$}; \node[above left,inner sep=1.5pt] at (H) {$H$};}
% Een teken voor een rechte hoek in Q, tussen QP en QR (ook in cavalièreperspectief): \rhoek{P}{Q}{R},
% of \rhoek[3mm]{P}{Q}{R} voor een groter teken
\newcommand{\rhoek}[4][2mm]{\draw[grijs,thin] ($(#3)!#1!(#2)$) -- ($(#3)!#1!(#2) + (#3)!#1!(#4) - (#3)$) -- ($(#3)!#1!(#4)$);}
% Een rechte door twee benoemde punten, aan beide kanten verlengd: \rechte[opties]{A}{B}
\newcommand{\rechte}[3][]{\draw[accent,thick,shorten <=-4mm,shorten >=-4mm,#1] (#2) -- (#3);}
% Een piramide T.ABCD met een vierkant grondvlak (zijde s) en hoogte h, de top boven het midden:
% \draadpiramide{s}{h}; zet (A), (B), (C), (D), (M) (midden van het grondvlak) en (T). D is verborgen.
\newcommand{\draadpiramide}[2]{%
  \coordinate (A) at (0,0,0); \coordinate (B) at (#1,0,0); \coordinate (C) at (#1,0,#1); \coordinate (D) at (0,0,#1);
  \coordinate (M) at (#1/2,0,#1/2); \coordinate (T) at (#1/2,#2,#1/2);
  \fill[accent!5] (A) -- (B) -- (C) -- (T) -- cycle;
  \draw[verborgen] (D) -- (A) (D) -- (C) (D) -- (T);
  \draw[ribbe] (A) -- (B) -- (C) -- (T) -- cycle (B) -- (T);}
\newcommand{\piramidenamen}{%
  \node[below left,inner sep=1.5pt] at (A) {$A$}; \node[below right,inner sep=1.5pt] at (B) {$B$};
  \node[right,inner sep=2pt] at (C) {$C$}; \node[left,inner sep=2pt] at (D) {$D$};
  \node[above,inner sep=2pt] at (T) {$T$};}
% Een vlak als parallellogram, horizontaal getekend: \vlakje{x}{y}{breedte}{diepte}{naam}
% (gewone, platte coördinaten; de naam staat in de hoek vooraan links)
\newcommand{\vlakje}[5]{%
  \filldraw[vlakvulling] (#1,#2) -- ++(#3,0) -- ++(#4*.7,#4*.7) -- ++(-#3,0) -- cycle;
  \node[accent,font=\small] at ($(#1,#2)+(.45,.22)$) {$#5$};}
